% Topic 1.1.13 — Fourier Transform.

\section{Преобразование Фурье}
\label{sec:1.1.13-fourier-transform}

\begin{ticket}
Преобразование Фурье абсолютно интегрируемой функции и его свойства. Преобразование Фурье производной и производная преобразования Фурье.
\end{ticket}

\subsection*{Теория}

Работаем на вещественной прямой, используя пространство
\[
  L^1(\R) \;=\; \Bigl\{ f \colon \R \to \C \;:\; f \text{ измерима},\; \|f\|_1 := \int_{\R} |f(x)|\, dx < \infty \Bigr\}.
\]
Функцию из $L^1(\R)$ называют \emph{абсолютно интегрируемой}.
Подробное изложение всех результатов раздела~--- в
\cite[Т.~II, гл.~XVIII, \S~3]{Zorich1981} и
\cite[гл.~VIII]{KolmogorovFomin}.

\subsubsection*{Определение и базовая регулярность}

\begin{definition}[Преобразование Фурье]
\label{def:fourier-transform}
Для $f \in L^1(\R)$ \emph{преобразованием Фурье} функции~$f$
называется функция $\hat f \colon \R \to \C$, заданная равенством
\[
  \hat f(\xi) \;=\; \int_{\R} f(x)\, e^{-i \xi x}\, dx, \qquad \xi \in \R.
\]
Обозначаем также $\mathcal{F}\{f\}(\xi) = \hat f(\xi)$.
\end{definition}

\begin{remark}
В выбранной нормировке множителя $2\pi$ нет в прямом преобразовании;
эта константа возникает в формуле обращения
(\cref{thm:fourier-inversion}). В других учебниках множитель
расщепляют симметрично как $1/\sqrt{2\pi}$ перед обоими
преобразованиями.
\end{remark}

\begin{proposition}[Ограниченность]
\label{prop:ft-bounded}
Для всякой $f \in L^1(\R)$ и каждого $\xi \in \R$
$|\hat f(\xi)| \leq \|f\|_1$. В частности, $\hat f$ ограничено
числом~$\|f\|_1$.
\end{proposition}
\proofof{ft-bounded}

\begin{theorem}[Непрерывность преобразования Фурье]
\label{thm:ft-continuous}
Для всякой $f \in L^1(\R)$ функция $\hat f$ равномерно непрерывна
на~$\R$.
\end{theorem}
\proofof{ft-continuous}

\begin{theorem}[Лемма Римана--Лебега на $\R$]
\label{thm:rl-line}
Для всякой $f \in L^1(\R)$
\[
  \hat f(\xi) \;\to\; 0 \qquad \text{при } |\xi| \to \infty.
\]
\end{theorem}
\proofof{rl-line}

\subsubsection*{Алгебраические свойства}

\begin{proposition}[Линейность, сдвиг, модуляция, растяжение]
\label{prop:ft-algebra}
Пусть $f, g \in L^1(\R)$, $\alpha, \beta \in \C$,
$a \in \R \setminus \{0\}$, $h \in \R$. Тогда для каждого $\xi \in \R$
\begin{enumerate}[label=(\alph*)]
  \item \emph{линейность:} $\widehat{\alpha f + \beta g}(\xi) = \alpha \hat f(\xi) + \beta \hat g(\xi)$;
  \item \emph{сдвиг по аргументу:} если $f_h(x) = f(x - h)$, то $\widehat{f_h}(\xi) = e^{-i \xi h} \hat f(\xi)$;
  \item \emph{модуляция:} если $g(x) = e^{i \omega x} f(x)$ при некотором $\omega \in \R$, то $\hat g(\xi) = \hat f(\xi - \omega)$;
  \item \emph{растяжение:} если $g(x) = f(a x)$, то $\hat g(\xi) = \tfrac{1}{|a|}\, \hat f\!\bigl(\tfrac{\xi}{a}\bigr)$.
\end{enumerate}
\end{proposition}
\proofof{ft-algebra}

\subsubsection*{Правила дифференцирования}

Две ключевые для этого билета формулы переводят гладкость в убывание и
обратно: дифференцирование~$f$ соответствует умножению $\hat f$ на
$i\xi$, а умножение~$f$ на~$x$~--- дифференцированию $\hat f$.

\begin{theorem}[Преобразование Фурье производной]
\label{thm:ft-of-derivative}
Пусть $f \in L^1(\R)$ всюду дифференцируема, $f' \in L^1(\R)$ и
$f(x) \to 0$ при $|x| \to \infty$. Тогда
\[
  \widehat{f'}(\xi) \;=\; i \xi \, \hat f(\xi), \qquad \xi \in \R.
\]
В общем случае: если~$f$ дифференцируема $k$ раз,
$f^{(j)} \in L^1(\R)$ при $j = 0, 1, \dots, k$ и
$f^{(j)}(x) \to 0$ при $|x| \to \infty$ для $j = 0, \dots, k - 1$, то
$\widehat{f^{(k)}}(\xi) = (i\xi)^k \, \hat f(\xi)$. В частности,
$|\xi|^k\, |\hat f(\xi)| \leq \|f^{(k)}\|_1$, то есть более гладкая
$f$ даёт быстрее убывающее $\hat f$.
\end{theorem}
\proofof{ft-of-derivative}

\begin{theorem}[Производная преобразования Фурье]
\label{thm:derivative-of-ft}
Пусть $f \in L^1(\R)$ и $x \mapsto x\, f(x) \in L^1(\R)$. Тогда
$\hat f$ непрерывно дифференцируема на~$\R$ и
\[
  (\hat f)'(\xi) \;=\; \widehat{(- i x\, f)}(\xi) \;=\; -i \int_{\R} x\, f(x)\, e^{-i \xi x}\, dx.
\]
В общем случае: если $x^j f(x) \in L^1(\R)$ при $j = 0, 1, \dots, k$,
то $\hat f \in C^k(\R)$ и
$(\hat f)^{(k)}(\xi) = \widehat{(-i x)^k f}(\xi)$. Более быстро
убывающая~$f$ даёт более гладкое~$\hat f$.
\end{theorem}
\proofof{derivative-of-ft}

\subsubsection*{Формула обращения (формулировка)}

\begin{theorem}[Формула обращения Фурье]
\label{thm:fourier-inversion}
Пусть $f \in L^1(\R)$ и $\hat f \in L^1(\R)$. Тогда для почти каждого
$x$
\[
  f(x) \;=\; \frac{1}{2\pi} \int_{\R} \hat f(\xi)\, e^{i \xi x}\, d\xi,
\]
причём равенство выполнено в каждой точке непрерывности~$f$. В этом
билете доказательство не приводим~--- см.\
\cite[Т.~II, гл.~XVIII, \S~3]{Zorich1981}; теорему используем лишь для
интерпретации примеров.
\end{theorem}

\subsection*{Доказательства}

\begin{proof}[\Cref{prop:ft-bounded}]
\proofanchor{ft-bounded}
Применим неравенство треугольника к интегралу:
\[
  |\hat f(\xi)| \;=\; \Bigl| \int_{\R} f(x)\, e^{-i \xi x}\, dx \Bigr|
  \;\leq\; \int_{\R} |f(x)|\, |e^{-i \xi x}|\, dx
  \;=\; \int_{\R} |f(x)|\, dx
  \;=\; \|f\|_1.
\]
\end{proof}

\begin{proof}[\Cref{thm:ft-continuous}]
\proofanchor{ft-continuous}
Для $\xi, \eta \in \R$
\[
  \hat f(\xi) - \hat f(\eta)
  \;=\; \int_{\R} f(x)\, \bigl(e^{-i \xi x} - e^{-i \eta x}\bigr)\, dx,
\]
откуда
\[
  |\hat f(\xi) - \hat f(\eta)|
  \;\leq\; \int_{\R} |f(x)|\, |e^{-i \xi x} - e^{-i \eta x}|\, dx.
\]
Подынтегральная функция мажорируется $2 |f(x)| \in L^1(\R)$ и
поточечно стремится к нулю при $|\xi - \eta| \to 0$. По теореме
Лебега о мажорированной сходимости
$|\hat f(\xi) - \hat f(\eta)| \to 0$ при $|\xi - \eta| \to 0$, причём
оценка зависит от $|\xi - \eta|$ только через подынтегральную
мажоранту,~--- значит, сходимость равномерна по~$\xi$.
\end{proof}

\begin{proof}[\Cref{thm:rl-line}]
\proofanchor{rl-line}
Сначала возьмём $f = \mathbf{1}_{[a, b]}$. Прямое вычисление даёт
\[
  \hat f(\xi) = \int_a^b e^{-i \xi x}\, dx = \frac{e^{-i \xi a} - e^{-i \xi b}}{i \xi}, \qquad \xi \neq 0,
\]
так что $|\hat f(\xi)| \leq 2/|\xi| \to 0$ при $|\xi| \to \infty$. По
линейности то же выполнено для всякой ступенчатой функции с конечным
числом значений на ограниченном интервале.

Для произвольной $f \in L^1(\R)$ и $\varepsilon > 0$ выберем
ступенчатую функцию~$\varphi$ с компактным носителем, для которой
$\|f - \varphi\|_1 < \varepsilon/2$,~--- это возможно благодаря
плотности ступенчатых функций с компактным носителем в $L^1(\R)$. По
\cref{prop:ft-bounded}
$|\hat f(\xi) - \hat \varphi(\xi)| \leq \|f - \varphi\|_1 < \varepsilon/2$.
Возьмём $\Lambda > 0$ так, чтобы $|\hat \varphi(\xi)| < \varepsilon/2$
при $|\xi| \geq \Lambda$. Тогда при таких~$\xi$
$|\hat f(\xi)| \leq |\hat \varphi(\xi)| + \varepsilon/2 < \varepsilon$.
\end{proof}

\begin{proof}[\Cref{prop:ft-algebra}]
\proofanchor{ft-algebra}
\emph{(а)} Линейность интеграла.

\emph{(б)} Замена $y = x - h$, $dy = dx$:
\[
  \widehat{f_h}(\xi) = \int_{\R} f(x - h)\, e^{-i \xi x}\, dx
                     = \int_{\R} f(y)\, e^{-i \xi (y + h)}\, dy
                     = e^{-i \xi h}\, \hat f(\xi).
\]

\emph{(в)} Прямой подсчёт:
$\hat g(\xi) = \int e^{i \omega x} f(x) e^{-i \xi x}\, dx
= \int f(x) e^{-i (\xi - \omega) x}\, dx = \hat f(\xi - \omega)$.

\emph{(г)} Замена $y = a x$, $dy = a\, dx$. При $a > 0$
\[
  \hat g(\xi) = \int_{\R} f(a x) e^{-i \xi x}\, dx
              = \frac{1}{a} \int_{\R} f(y) e^{-i (\xi/a) y}\, dy
              = \frac{1}{a}\, \hat f(\xi/a).
\]
При $a < 0$ замена меняет местами пределы интегрирования, и получается
$\hat g(\xi) = -\tfrac{1}{a}\, \hat f(\xi/a) = \tfrac{1}{|a|}\, \hat f(\xi/a)$.
\end{proof}

\begin{proof}[\Cref{thm:ft-of-derivative}]
\proofanchor{ft-of-derivative}
Для $R > 0$ применим интегрирование по частям на $[-R, R]$ при
$u = e^{-i\xi x}$, $dv = f'(x)\, dx$:
\[
  \int_{-R}^{R} f'(x) e^{-i\xi x}\, dx
  = \bigl[f(x) e^{-i\xi x}\bigr]_{-R}^{R}
    + i\xi \int_{-R}^{R} f(x) e^{-i\xi x}\, dx.
\]
Граничный член равен $f(R) e^{-i\xi R} - f(-R) e^{i\xi R}$, его
модуль не превосходит $|f(R)| + |f(-R)|$, а это по условию стремится
к нулю при $R \to \infty$. Оба интеграла стремятся к
$\widehat{f'}(\xi)$ и $\hat f(\xi)$ соответственно (поскольку
$f, f' \in L^1$). Переходя к пределу при $R \to \infty$, получаем
$\widehat{f'}(\xi) = i\xi\, \hat f(\xi)$. По индукции
$\widehat{f^{(k)}}(\xi) = (i\xi)^k \hat f(\xi)$, и
\cref{prop:ft-bounded}, применённая к~$f^{(k)}$, даёт
$|\xi|^k |\hat f(\xi)| = |\widehat{f^{(k)}}(\xi)| \leq \|f^{(k)}\|_1$.
\end{proof}

\begin{proof}[\Cref{thm:derivative-of-ft}]
\proofanchor{derivative-of-ft}
Положим $F(\xi, x) = f(x) e^{-i \xi x}$, тогда
$\hat f(\xi) = \int_{\R} F(\xi, x)\, dx$. Частная производная
$\partial_\xi F(\xi, x) = -i x\, f(x)\, e^{-i \xi x}$ существует при
любых $(\xi, x)$ и мажорируется функцией $|x f(x)|$, лежащей в
$L^1(\R)$ по условию. Применима формула Лейбница для дифференцирования
интеграла, зависящего от параметра, см.\
\cite[Т.~II, гл.~XVII]{Zorich1981}, и она даёт
\[
  (\hat f)'(\xi) \;=\; \int_{\R} \partial_\xi F(\xi, x)\, dx
                 \;=\; -i \int_{\R} x\, f(x)\, e^{-i \xi x}\, dx
                 \;=\; \widehat{(- i x\, f)}(\xi).
\]
Правая часть~--- преобразование Фурье $L^1$-функции, а значит,
непрерывна по~$\xi$ согласно \cref{thm:ft-continuous}; следовательно,
$\hat f \in C^1(\R)$. Общий случай $C^k$ доказывается по индукции при
выполнении условия $x^j f(x) \in L^1$.
\end{proof}

\subsection*{Примеры}

\begin{example}[Индикатор интервала]
\label{ex:ft-indicator}
Пусть $f = \mathbf{1}_{[-1, 1]}$. Тогда $f \in L^1$, и при $\xi \neq 0$
\[
  \hat f(\xi) = \int_{-1}^{1} e^{-i \xi x}\, dx
              = \frac{e^{-i \xi} - e^{i \xi}}{-i\xi}
              = \frac{2 \sin \xi}{\xi}.
\]
В точке $\xi = 0$ имеем $\hat f(0) = 2 = \lim_{\xi \to 0} (2 \sin \xi)/\xi$.
Преобразование вещественно, чётно, непрерывно и стремится к нулю на
бесконечности~--- что согласуется с
\cref{thm:ft-continuous,thm:rl-line}. Заметим, однако, что
$\hat f \notin L^1(\R)$, поскольку $|2 \sin \xi / \xi|$ не абсолютно
интегрируемо (осциллирующий хвост порядка $1/\xi$); поэтому формулу
обращения нужно применять в подходящем смысле.
\end{example}

\begin{example}[Двусторонняя экспонента]
\label{ex:ft-double-exp}
Пусть $f(x) = e^{-|x|}$. Тогда $f \in L^1$, $\|f\|_1 = 2$ и
\begin{align*}
  \hat f(\xi)
  &= \int_{-\infty}^{0} e^{x}\, e^{-i\xi x}\, dx + \int_{0}^{\infty} e^{-x}\, e^{-i\xi x}\, dx \\
  &= \frac{1}{1 - i\xi} + \frac{1}{1 + i\xi}
   = \frac{2}{1 + \xi^2}.
\end{align*}
И $f$, и $\hat f$ интегрируемы; формула обращения восстанавливает~$f$
из $\hat f$ поточечно. Преобразование убывает как $\xi^{-2}$~---
согласуется с тем, что~$f$ принадлежит лишь классу $C^0$ (одна
обобщённая производная в~$L^1$ даёт дополнительный множитель $1/\xi$ в
\cref{thm:ft-of-derivative}).
\end{example}

\begin{example}[Гауссиана: самосопряжённая функция]
\label{ex:ft-gaussian}
Пусть $f(x) = e^{-x^2/2}$. Тогда $f \in L^1$ и можно показать, что
$\hat f(\xi) = \sqrt{2\pi}\, e^{-\xi^2/2}$. Изящное доказательство
опирается на правила дифференцирования: $f'(x) = -x f(x)$, поэтому,
применяя $\mathcal{F}$ к обеим частям и пользуясь
\cref{thm:ft-of-derivative,thm:derivative-of-ft},
\[
  i \xi\, \hat f(\xi) \;=\; \widehat{f'}(\xi) \;=\; \widehat{-x f}(\xi) \;=\; -i\, (\hat f)'(\xi),
\]
то есть $(\hat f)'(\xi) = -\xi\, \hat f(\xi)$. Это то же
дифференциальное уравнение первого порядка, которому удовлетворяет
сама~$f$ (с заменой $x$ на $\xi$); следовательно,
$\hat f(\xi) = C\, e^{-\xi^2/2}$ при некоторой константе~$C$.
В точке $\xi = 0$:
$C = \hat f(0) = \int_{\R} e^{-x^2/2}\, dx = \sqrt{2\pi}$. Итог:
$\hat f(\xi) = \sqrt{2\pi}\, e^{-\xi^2/2}$~--- гауссиана с точностью
до постоянного множителя является собственным преобразованием Фурье.
\end{example}

\begin{problem}[Соответствие гладкость~--- убывание]
\label{prob:smoothness-decay}
Пусть $f \in L^1(\R)$, $f \in C^k(\R)$, $f, f', \dots, f^{(k)} \in L^1(\R)$
и $f^{(j)}(x) \to 0$ при $|x| \to \infty$ для $j = 0, \dots, k - 1$.
Докажите, что $\hat f(\xi) = o(|\xi|^{-k})$ при $|\xi| \to \infty$.
(Указание: воспользуйтесь \cref{thm:ft-of-derivative} вместе с
\cref{thm:rl-line}.)
\end{problem}

\begin{proof}[Решение]
Применяя \cref{thm:ft-of-derivative} к производной порядка~$k$
(условия теоремы выполнены ровно по предположениям задачи), получаем
$\widehat{f^{(k)}}(\xi) = (i\xi)^k\, \hat f(\xi)$, откуда
$|\xi|^k\, |\hat f(\xi)| = \bigl|\widehat{f^{(k)}}(\xi)\bigr|$. Так как
$f^{(k)} \in L^1(\R)$, по \cref{thm:rl-line}
$\widehat{f^{(k)}}(\xi) \to 0$ при $|\xi| \to \infty$, то есть
$|\xi|^k\, |\hat f(\xi)| \to 0$. Это и есть $\hat f(\xi) = o(|\xi|^{-k})$.
\end{proof}
